\section{Detect a Loop in a Linked List}
\label{sec:ll-detect-loop}

\problemstatement{Given the head of a singly linked list, determine whether
it contains a cycle (some node's \texttt{next} points back to an earlier
node), using $O(1)$ extra space.}

\begin{patternbox}
\emph{``Is there a cycle?''} is \textbf{Floyd's tortoise--hare}: a slow
pointer (one step) and a fast pointer (two steps). If there is a loop, the
fast pointer laps the slow one and they meet inside the loop; if there is no
loop, the fast pointer reaches null.
\end{patternbox}

\subsection*{Why the two pointers must meet}

Inside a loop, each iteration the fast pointer gains one node on the slow
pointer, so the gap shrinks by one every step and eventually hits zero --
they collide. If the list is acyclic, the fast pointer falls off the end
(reaches null) first. This detects a cycle in one pass with only two
pointers, unlike a hash-set approach that needs $O(n)$ memory.

\begin{ideabox}[Gaining One Node per Step Guarantees a Meeting]
Relative to the slow pointer, the fast pointer moves one node closer each
iteration. On a finite cycle the distance between them is bounded, so it
must reach zero -- a collision -- proving a loop exists.
\end{ideabox}

\subsection*{Algorithm}

\begin{enumerate}
  \item \texttt{slow = fast = head}.
  \item While \texttt{fast} and \texttt{fast->next}: advance \texttt{slow}
        by one, \texttt{fast} by two; if they are equal, return true.
  \item If the loop exits, return false.
\end{enumerate}

\subsection*{Dry run}

\dryrun{$1\to2\to3\to4\to2$ (node $4$ links back to $2$).}

\begin{itemize}
  \item slow/fast advance; fast laps and meets slow at some node inside
        $\{2,3,4\}$.
  \item Meeting confirms a cycle -- return true.
\end{itemize}

\subsection*{C++ implementation}

\begin{lstlisting}
#include <bits/stdc++.h>
using namespace std;

struct ListNode {
    int val; ListNode* next;
    ListNode(int v) : val(v), next(nullptr) {}
};

bool hasCycle(ListNode* head) {
    ListNode* slow = head;
    ListNode* fast = head;
    while (fast && fast->next) {
        slow = slow->next;
        fast = fast->next->next;
        if (slow == fast) return true;            // collision inside loop
    }
    return false;
}

int main() {
    ListNode* a = new ListNode(1);
    a->next = new ListNode(2);
    a->next->next = new ListNode(3);
    a->next->next->next = new ListNode(4);
    a->next->next->next->next = a->next;          // 4 -> 2, a cycle
    cout << (hasCycle(a) ? "cycle" : "no cycle") << "\n";
    return 0;
}
\end{lstlisting}

\begin{complexitybox}
\textbf{Time Complexity.} $O(n)$. \textbf{Space Complexity.} $O(1)$ (a hash
set would cost $O(n)$).
\end{complexitybox}

\begin{takeawaybox}
Floyd's cycle detection uses two speeds and constant memory; a collision
proves a loop. The same meeting point locates the loop's start and length
(next two sections), so this is the gateway to all cycle problems.
\end{takeawaybox}
